\clearpage
\subsection{Mengenoperationen}
\label{subsec:fundamentals_set_operations}

Zwischen Mengen existieren Operationen, mit denen aus gegebenen Mengen neue
Mengen gebildet werden können. Zu den wichtigsten gehören Schnitt, Vereinigung
und Differenz.

\subsubsection{Schnittmenge}
\label{subsubsec:fundamentals_set_intersection}

\noindent\begin{minipage}{\linewidth}
\begin{DefinitionBox}{Schnittmenge}{fundamentals_set_intersection}
  \begin{minipage}[c]{0.70\linewidth}
    Die \emph{Schnittmenge} \EnglishTerm{intersection} zweier Mengen $A$ und
    $B$ enthält alle Elemente, die sowohl in der Menge $A$ als auch in der
    Menge $B$ enthalten sind.
    \begin{ClozeArea}[24]{3}
      \[
        A\cap B=\{x\mid x\in A\land x\in B\}.
      \]
    \end{ClozeArea}
  \end{minipage}\hfill
  \begin{minipage}[c]{0.25\linewidth}
    \centering
    \begin{tikzpicture}[scale=1.2]
  \begin{scope}
    \clip (-0.5,0) circle (0.75);
    \fill[THMBlue!30] (0.5,0) circle (0.75);
  \end{scope}
  \draw[thick] (-0.5,0) circle (0.75);
  \draw[thick] (0.5,0) circle (0.75);
  \node at (-0.5,0) {$A$};
  \node at (0.5,0) {$B$};
\end{tikzpicture}

  \end{minipage}
\end{DefinitionBox}

\begin{ExampleBox}{Schnittmenge}{fundamentals_set_intersection_example}
  \begin{enumerate}[label=(\alph*),itemsep=0.6em]
    \item $\{1,2\}\cap\{2,3\}=\Cloze{3cm}{\{2\}}$
    \item $\{1,2,3\}\cap\{1,2,3\}=\Cloze{3cm}{\{1,2,3\}}$
    \item $M\cap\emptyset=\Cloze{2.5cm}{\emptyset}\quad(\forall M)$
  \end{enumerate}
\end{ExampleBox}
\end{minipage}



\subsubsection{Vereinigungsmenge}
\label{subsubsec:fundamentals_set_union}

\noindent\begin{minipage}{\linewidth}
\begin{DefinitionBox}{Vereinigungsmenge}{fundamentals_set_union}
  \begin{minipage}[c]{0.70\linewidth}
    Die \emph{Vereinigungsmenge} \EnglishTerm{union} zweier Mengen $A$ und
    $B$ enthält alle Elemente, die in der Menge $A$, in der Menge $B$ oder in
    beiden Mengen enthalten sind.
    \begin{ClozeArea}[24]{3}
      \[
        A\cup B=\{x\mid x\in A\lor x\in B\}.
      \]
    \end{ClozeArea}
  \end{minipage}\hfill
  \begin{minipage}[c]{0.25\linewidth}
    \centering
    \begin{tikzpicture}[scale=1.2]
  \fill[THMBlue!30] (-0.5,0) circle (0.75);
  \fill[THMBlue!30] (0.5,0) circle (0.75);
  \draw[thick] (-0.5,0) circle (0.75);
  \draw[thick] (0.5,0) circle (0.75);
  \node at (-0.5,0) {$A$};
  \node at (0.5,0) {$B$};
\end{tikzpicture}

  \end{minipage}
\end{DefinitionBox}

\begin{ExampleBox}{Vereinigungsmenge}{fundamentals_set_union_example}
  \begin{enumerate}[label=(\alph*),itemsep=0.6em]
    \item $\{1,2\}\cup\{2,3\}=\Cloze{3cm}{\{1,2,3\}}$
    \item $\{1,2,3\}\cup\{1,2,3\}=\Cloze{3cm}{\{1,2,3\}}$
    \item $M\cup\emptyset=\Cloze{2.5cm}{M}\quad(\forall M)$
  \end{enumerate}
\end{ExampleBox}
\end{minipage}


\subsubsection{Differenzmenge}
\label{subsubsec:fundamentals_set_difference}

\noindent\begin{minipage}{\linewidth}
\begin{DefinitionBox}{Differenzmenge}{fundamentals_set_difference}
  \begin{minipage}[c]{0.70\linewidth}
    Die \emph{Differenzmenge} \EnglishTerm{set difference} $A\setminus B$
    enthält alle Elemente, die in der Menge $A$, aber nicht in der Menge $B$
    enthalten sind.
    \begin{ClozeArea}[24]{3}
      \[
        A\setminus B=\{x\mid x\in A\land x\notin B\}.
      \]
    \end{ClozeArea}
  \end{minipage}\hfill
  \begin{minipage}[c]{0.25\linewidth}
    \centering
    \begin{tikzpicture}[scale=1.2]
  \begin{scope}
    \clip (-0.5,0) circle (0.75);
    \fill[THMBlue!30] (-0.5,0) circle (0.75);
    \clip (0.5,0) circle (0.75);
    \fill[white] (0.5,0) circle (0.75);
  \end{scope}
  \draw[thick] (-0.5,0) circle (0.75);
  \draw[thick] (0.5,0) circle (0.75);
  \node at (-0.5,0) {$A$};
  \node at (0.5,0) {$B$};
\end{tikzpicture}

  \end{minipage}
\end{DefinitionBox}

\begin{ExampleBox}{Differenzmenge}{fundamentals_set_difference_example}
  \begin{enumerate}[label=(\alph*),itemsep=0.6em]
    \item $\{1,2\}\setminus\{2,3\}=\Cloze{3cm}{\{1\}}$
    \item $\{1,2,3\}\setminus\{1,2,3\}=\Cloze{2.5cm}{\emptyset}$
    \item $M \setminus \emptyset = \Cloze{2.5cm}{M}$
  \end{enumerate}
\end{ExampleBox}
\end{minipage}
